\section{Methods}

\subsection{Testbed and Input}
The person remains seated while controlling a virtual climbing task built
in Unity using OpenXR~\cite{openxr} and the XR Interaction
Toolkit~\cite{unity-xri}, connected through Meta Quest Link. The environment
has no falling or momentum; upward movement is capped at 0.25\,m/s, and a
visual vignette is applied~\cite{fernandes2016combating}.

Mouse movement positions the active hand along the wall. Pressing the grab
button attaches the hand to a reachable hold; dragging downward then raises
the virtual body. Releasing the button leaves that hand attached and switches
control to the other hand. An unsuccessful grab preserves existing support.
The grab action uses the left mouse button, with Space used in earlier runs.
E/Q provide manual depth adjustment at 0.15\,m/s.

\subsection{Simultaneous Control and Authority}
The division of control is fixed: the person controls lateral hand movement,
grabbing, and climbing, while assistance adjusts reaching depth. Unity supplies
the Python service with hand positions, the active hand, reaching and gripping
states, and hold geometry. The service uses this information to infer a target
hold and calculate bounded depth adjustments. Images, language, and experimental
target cues are not used for inference.

During each Unity update, human-controlled lateral coordinates $(x,y)$ are
combined with the accepted assisted depth coordinate $z$ to form one hand pose
(Fig.~\ref{fig:arch}). The person can continue moving the mouse while assistance
updates depth. Both therefore contribute during the same reach, even though
the service and Unity update asynchronously.

The person retains control over grabbing and depth override. Grabbing freezes
depth, and pressing E/Q transfers depth control to the person for the remainder
of that reach. Unity rejects stale or invalid assistance updates and handles
interruptions locally; resuming after an interruption requires deliberate
user action.

\subsection{Depth Target and Conditions}
Hand contact uses a 0.04\,m-radius sphere; holds protrude 0.04, 0.08, or
0.12\,m. Assistance targets near-surface contact within a hold's lateral grasp
region, where valid depth exists. Grabs accept a broader overlap interval;
retained depth can permit the next grab without adjustment.

Conditions are \emph{manual} (M; E/Q), \emph{assist-nearest}
(N; the nearest unoccupied hold whose lateral grasp region contains the hand),
and \emph{assist-intent} (I; motion-inferred target). Both assisted conditions
share a depth controller; intent adds the selection rules below.

\begin{figure}[!t]
\centering
% Fig. 2: conditions for moving depth; exact thresholds remain in Methods.
\begin{tikzpicture}[
  font=\footnotesize, >={Stealth[length=4pt]},
  box/.style={draw, rounded corners=2pt, align=center, inner sep=4pt,
    text width=4.50cm},
  flow/.style={->, thick}
]
\node[box, fill=orange!10] (infer) at (0,0)
  {\textbf{1. Predict the intended hold}\\Project recent lateral hand motion};
\node[box, fill=orange!10] (check) at (0,-1.35)
  {\textbf{2. Check before moving}\\Confident + stable target?\\Inside grasp region + valid depth?};
\node[box, fill=orange!10] (move) at (0,-2.75)
  {\textbf{3. Advance depth toward contact}\\Small steps coupled to lateral motion};
\node[draw, rounded corners=2pt, fill=gray!7,
  align=center, text width=1.60cm, inner sep=4pt] (wait) at (3.95,-1.35)
  {\textbf{Hold depth}\\Keep observing};
\draw[flow] (infer) -- (check);
\draw[flow] (check) -- node[right, fill=white, inner sep=1pt] {yes} (move);
\draw[flow] (check) -- node[above, inner sep=1pt] {no} (wait);
\end{tikzpicture}

\caption{Intent-assisted depth (nearest has no confidence gate). Uncertainty holds
depth. After lateral motion stops, the depth budget ends after 150\,ms.}
\label{fig:gate}
\end{figure}

\subsection{Intent Inference and Target Selection}
At 10\,Hz, intent estimates lateral velocity $\vec{v}$ over 250\,ms,
caps speed at 0.3\,m/s, and projects lateral position
$\vec{q}$ to $\hat{\vec{q}}=\vec{q}+(0.5\,\mathrm{s})\vec{v}$.
For unoccupied hold $i$ with lateral centre $\vec{c}_i$ and grasp radius
$r_i$, predicted distance outside its grasp region is
$d_i=\max(0,\|\hat{\vec{q}}-\vec{c}_i\|-r_i)$.
Weights $w_i=\exp(-d_i^2/2\sigma^2)$, with $\sigma=0.08$\,m, and a
``none'' weight of 0.1 are normalized to scores $p_i$.

Selection requires score $p_i\ge0.80$, a 0.25 lead over the runner-up,
and valid grasp geometry and depth, maintained for 200\,ms. A selected target's
score below 0.55 holds depth; persistence for 200\,ms withdraws the target.
Invalid geometry withdraws it immediately (Fig.~\ref{fig:gate}). These rules
follow prediction and arbitration~\cite{dragan2013policy,javdani2015shared};
scores are uncalibrated, and thresholds and hysteresis are design choices.

\subsection{Motion-Coupled Depth Control}
Motion coupling limits both assisted conditions using lateral hand speed.
For eligible target depth $z^*$ and current depth $z_t$, the maximum step
$b_t$ and commanded change are
\begin{equation}
\begin{aligned}
b_t &= \min(0.015\,\mathrm{m},\; 0.6\,\mathrm{m/s}\cdot\Delta t,\; k\tilde v_t\Delta t),\\
\Delta z_t &= \operatorname{sgn}(z^*-z_t)\min(|z^*-z_t|,b_t).
\end{aligned}
\end{equation}
Here $\Delta t$ is the observation interval and $k=1.5$. Effective speed
$\tilde v_t$ uses lateral speed at or above 1\,cm/s, retaining the last such
speed for a 0.15\,s grace period, then zero. Assistance thus continues briefly
after stopping, then holds depth. Steps never exceed
1.5\,cm, within the client's 2\,cm limit. Earlier coupled runs used $k=0.75$
and a 0.15\,m/s cap.

\subsection{Evaluation Design and Measures}
Table~\ref{tab:evidence} summarizes simulations, integration trials, and headset
sessions examining controller tradeoffs, signal delivery, and depth effort,
respectively. The intent sweep served both tuning and evaluation; results
are descriptive.

One developer tested two headset sessions. Each planned one practice
trial/condition and six three-trial blocks: M-N-I-I-N-M, then N-I-M-M-I-N.
Trials required four cued holds; only completed measured trials enter separate
session summaries. Block depth effort used ``Getting the depth right took
effort'' (1: none; 7: a lot). Discomfort $\ge4$ or no rating within 60\,s
stopped the session.

Trial time spans ready to finish. Assistance-state episodes span entry to a
state/target change and are scored against each overlapping reach/cue interval;
wrong-target episodes select a non-cued hold. Headset \emph{engagement coverage}
is the fraction of intervals containing an episode, summarized by the trial
median. States need not produce movement. Wrong-hold travel measures applied
depth changes toward non-cued holds. Simulation coverage averages within-reach
time assisted toward the cue. Causal checks flag unauthorized movement or
direction/cap violations.

Simulated \emph{extra motion} is lateral movement added after lateral arrival
and a depth stall of about 0.3\,s. The headset metric requires a completed cued
reach with a logged coupling hold of at least 0.3\,s followed by resumed
assistance.
